% =============================================================================
% TRABALHO FINAL: ESTRUTURAS EM ÁRVORES E TEORIA DOS GRAFOS
% Disciplina: Algoritmos e Estruturas de Dados II (AEDS II)
% Orientador: Prof. Michel Pires da Silva (CEFET-MG / UFCI-MG)
% Modelo: IEEE Conference Template
% =============================================================================

\documentclass[conference]{IEEEtran}

% --- Pacotes Básicos e Idioma ---
\usepackage[utf8]{inputenc}
\usepackage[T1]{fontenc}
\usepackage[brazilian]{babel}

% --- Pacotes Matemáticos e Algoritmos ---
\usepackage{amsmath,amssymb,amsfonts}
\usepackage{algorithmic}

% --- Pacotes Gráficos e Tabelas ---
\usepackage{graphicx}
\usepackage{booktabs}
\usepackage{xcolor}
\usepackage{textcomp}

% --- Citações e Links ---
\usepackage{cite}
\usepackage{url}

\begin{document}

% =============================================================================
% TÍTULO DO ARTIGO
% =============================================================================
\title{Análise de Contratações Públicas Através da Teoria dos Grafos: congluios, monopólios e oligopólios}

% =============================================================================
% AUTORIA (Regras da Seção 3.1 do edital):
% 1. Primeiros autores: Alunos da equipe (até 3 integrantes), ordenados por
%    nível de contribuição (ou ordem alfabética se contribuição igualitária).
% 2. Último autor: OBRIGATORIAMENTE o professor orientador (Michel Pires da Silva).
% =============================================================================
\author{
    \IEEEauthorblockN{Igor Peixoto Rodrigues}
    \IEEEauthorblockA{\textit{Departamento de Computação} \\
    \textit{UFCI-MG} \\
    Divinópolis, Brasil \\
    igorpeixoto@cefetmg.br.com}
    \and
    \IEEEauthorblockN{Pedro Henrique de Araújo Carvalho}
    \IEEEauthorblockA{\textit{Departamento de Computação} \\
    \textit{UFCI-MG} \\
    Divinópolis, Brasil \\
    ypedroaraujo@gmail.com}
    \and
    \IEEEauthorblockN{Michel Pires da Silva}
    \IEEEauthorblockA{\textit{Departamento de Computação} \\
    \textit{UFCI-MG} \\
    Divinópolis, Brasil \\
    michel@cefetmg.br}
}

\maketitle

% =============================================================================
% ESTRUTURA MODULAR DO ARTIGO (SEÇÕES INDIVIDUAIS EM report/sections/)
% =============================================================================

% Resumo e Palavras-chave
% =============================================================================
% RESUMO / ABSTRACT
% Diretrizes do Professor (Seção 3.2 do edital):
% - Um parágrafo conciso (máximo de 250 palavras).
% - Deve responder diretamente às 4 perguntas:
%   1. Qual o problema?
%   2. Há soluções, quais?
%   3. Como foi modelado (qual estrutura de árvore/grafo)?
%   4. Qual o principal resultado obtido?
% =============================================================================

\begin{abstract}
% TODO: Escreva o resumo aqui (máximo de 250 palavras) respondendo diretamente
% às quatro perguntas exigidas no edital.
\end{abstract}

\begin{IEEEkeywords}
% TODO: Insira aqui as palavras-chave separadas por vírgula (3 a 6 termos).
% Exemplo: Teoria dos Grafos, Algoritmos e Estruturas de Dados, Redes Complexas.
\end{IEEEkeywords}


% 1. Introdução (Técnica do Funil, Relevância, Motivação, Objetivos e Roteiro)
\section{Introdução}
\label{sec:introducao}

A administração pública brasileira exerce papel central na prestação de serviços essenciais à população. É no âmbito municipal que se concretizam as políticas públicas cotidianas, abrangendo desde a atenção primária à saúde e a educação básica até a manutenção da malha viária urbana e os serviços de saneamento. A execução contínua dessas atribuições demanda gestão eficiente e mobilização de expressivos recursos públicos, viabilizados tanto por receitas tributárias próprias quanto por transferências constitucionais das esferas estadual e federal.

Para atender às demandas da coletividade, o poder público recorre rotineiramente ao mercado para aquisição de insumos médico-hospitalares, contratação de obras de engenharia civil, locação de frotas e prestação de serviços especializados. Dessa forma, consolida-se a interface entre o ente público contratante (prefeituras e secretarias municipais) e os fornecedores privados. Embora a legislação nacional — estabelecida pela Nova Lei de Licitações e Contratos Administrativos (Lei Federal nº 14.133/2021)~\cite{brasil2021licitacoes}, pela Lei de Defesa da Concorrência (Lei Federal nº 12.529/2011)~\cite{brasil2011concorrencia} e pela Lei Anticorrupção (Lei Federal nº 12.846/2013)~\cite{brasil2013anticorrupcao} — imponha rígidos princípios de publicidade, isonomia e ampla competitividade, práticas anticompetitivas e arranjos societários ou comerciais podem explorar brechas relacionais e temporais, dificultando sua detecção pelos mecanismos tradicionais de controle.

Essa dificuldade decorre do fato de que instrumentos convencionais de fiscalização operam sob perspectiva predominantemente tabular, avaliando certames, contratos e empenhos de maneira atômica e isolada. Sob essa visão fragmentada, verificações lineares têm baixa capacidade de capturar padrões sistêmicos complexos, tais como vínculos societários ou cadastrais ocultos entre empresas formalmente concorrentes, o rodízio coordenado de fornecedores vencedores com propostas de cobertura (\textit{cover bidding}) e a concentração orçamentária em poucos credores ou grupos econômicos. O sinal investigativo relevante, portanto, emerge do comportamento cumulativo e das relações estruturais consolidadas entre o poder público e o mercado fornecedor ao longo do tempo.

Para superar essa limitação, a Teoria dos Grafos fornece uma base analítica rigorosa. Ao modelar unidades orçamentárias, empresas licitantes e seus sócios reais como vértices ($V$) e ao representar participações, co-disputas, alternâncias de vitórias e fluxos empenhados como arestas ponderadas ($E$), a abordagem em redes transcende a inspeção isolada de tabelas. Essa perspectiva possibilita mapear conexões indiretas, isolar subgrupos suspeitos de co-atuação e quantificar assimetrias de dependência e concentração financeira que permaneceriam ocultas em auditorias convencionais.

A motivação prática para este estudo nasce da lacuna de granularidade observada nos portais locais de transparência: enquanto o portal municipal de Divinópolis/MG disponibiliza apenas visões agregadas~\cite{portal_transparencia_betha}, os dados jurisdicionados do Sistema Informatizado de Contas dos Municípios (SICOM), mantidos em dados abertos pelo Tribunal de Contas do Estado de Minas Gerais (TCE-MG)~\cite{tce_mg_dados_abertos,tcemg2025in01}, reúnem 55 tabelas relacionais que cobrem todo o ciclo da despesa pública. Assim, este trabalho tem como objetivo geral desenvolver e validar um \textit{pipeline} computacional baseado em Teoria dos Grafos para \textbf{triagem indiciária de risco} em licitações e contratos administrativos do município de Divinópolis (MG) no exercício consolidado de 2025. Em consonância com a jurisprudência do Tribunal de Contas da União (TCU)~\cite{tcu2011acordao2341,tcu2024acordao1798}, enfatiza-se que os indicadores estruturais obtidos não constituem prova automática de fraude ou ilícito, mas sim critérios quantitativos para priorização de auditoria.

De forma específica, a investigação estrutura-se em três eixos complementares: (i) no \textbf{Eixo Societário e de Materialidade}, construir um grafo heterogêneo multicamada cruzando co-disputas com o Quadro de Sócios e Administradores (QSA) da Receita Federal para pontuar vínculos ocultos, isolar triângulos societários ($K_3$) e ponderar o risco pelo valor efetivamente empenhado e pago; (ii) no \textbf{Eixo Comportamental}, modelar um dígrafo de disputas (Perdedor $\to$ Vencedor) para identificar inversões recíprocas de vitórias em datas distintas, ciclos curtos ($C_3/C_4$), comunidades de co-concorrência e supressão de desconto frente às cotações de referência; e (iii) no \textbf{Eixo de Concentração Orçamentária}, estruturar um grafo bipartido Unidade Orçamentária $\leftrightarrow$ Credor para mensurar o Índice Herfindahl--Hirschman (HHI), a entropia de dispersão de gastos, o \textit{PageRank} ponderado e o impacto de grupos societários ($\Delta\text{HHI}$), diferenciando a dependência estrutural de entes públicos ou filantrópicos da concentração em fornecedores privados.

O restante deste artigo está organizado da seguinte maneira. A Seção~\ref{sec:referencial} apresenta o referencial teórico, integrando o arcabouço jurídico-regulatório, as tipologias de risco em compras públicas e os algoritmos clássicos de grafos empregados. A Seção~\ref{sec:metodologia} detalha a arquitetura do \textit{pipeline} em dados do SICOM 2025, a modelagem formal de vértices e arestas, o protocolo de conformidade com a LGPD e o fluxograma completo da solução. A Seção~\ref{sec:resultados} expõe e discute criticamente os resultados empíricos dos três eixos, os achados consolidados, os gargalos computacionais e os tempos reais de execução. Por fim, a Seção~\ref{sec:conclusao} sintetiza as contribuições obtidas, analisa as limitações empíricas da base pública e aponta direções para trabalhos futuros. Com esta proposta, espera-se disponibilizar uma metodologia reprodutível, auditável e escalável capaz de apoiar órgãos de controle interno e externo na triagem preventiva de anomalias contratuais no setor público.


% 2. Referencial Teórico (Fundamentos de Grafos/Árvores e Estado da Arte)
\section{Referencial Teórico e Estado da Arte}
\label{sec:referencial}

A análise automatizada de riscos em contratações públicas articula fundamentos da Ciência da Computação (algoritmos em grafos e estruturas de dados), da Ciência das Redes Complexas e da regulação jurídica e antitruste das licitações. Esta seção consolida o arcabouço normativo, a literatura científica e as formulações matemáticas que sustentam os três eixos do estudo.

\subsection{Arcabouço Jurídico e Tipologias de Risco Licitatório}

As licitações públicas visam assegurar a isonomia entre competidores e selecionar a proposta mais vantajosa para a Administração, conforme os princípios expressos nos arts.~5º, 9º, inciso~I, alínea ``a'', e 11, incisos~I a III, da Lei nº 14.133/2021~\cite{brasil2021licitacoes}. O mesmo diploma impõe pesquisa de preços compatível com o mercado e consideração da economia de escala (art.~23, §§~1º e 4º), tipifica infrações administrativas graves (art.~155, incisos~VIII a XII), estabelece a dosimetria das sanções conforme o dano apurado (art.~156, §~1º, incisos~I e IV), autoriza a desconsideração da personalidade jurídica para alcançar sócios e empresas coligadas (art.~160) e inclui o art.~337-F no Código Penal para coibir a frustração do caráter competitivo.

Sob a ótica concorrencial, a Lei nº 12.529/2011 (art.~36, \textit{caput}, §~3º, incisos~I, alíneas ``c'' e ``d'', e II)~\cite{brasil2011concorrencia} enquadra como infrações à ordem econômica a divisão concertada de mercados e a combinação de preços ou abstenções em certames. Complementarmente, a Lei Anticorrupção (Lei nº 12.846/2013, art.~5º, inciso~IV)~\cite{brasil2013anticorrupcao} responsabiliza objetivamente pessoas jurídicas por atos lesivos em licitações, enquanto a Lei Complementar nº 123/2006 (art.~3º, §~4º, incisos~III a V)~\cite{brasil2006lc123} veda a fruição de tratamento diferenciado de microempresas (ME) e empresas de pequeno porte (EPP) quando há participação societária cruzada que exceda o teto legal. Já a Lei de Improbidade Administrativa (Lei nº 8.429/1992, com redação dada pela Lei nº 14.230/2021, art.~10 e art.~12, inciso~II)~\cite{brasil1992improbidade} exige demonstração de dolo específico e de perda patrimonial efetiva para responsabilização, o que fundamenta o cruzamento entre indícios estruturais e materialidade financeira. Por fim, a jurisprudência do TCU~\cite{tcu2011acordao2341,tcu2024acordao1798} consolida o entendimento de que a existência isolada de sócio em comum ou parentesco não configura ilícito por si só, exigindo a convergência de múltiplos indícios objetivos para justificar apuração aprofundada.

Na literatura econômica e nos guias da Organização para a Cooperação e Desenvolvimento Econômico (OCDE)~\cite{oecd2025guidelines} e do Conselho Administrativo de Defesa Econômica (CADE)~\cite{cade2019guia}, os arranjos restritivos à competição em compras públicas manifestam-se em quatro mecanismos principais: (i) \textit{cover bidding} (propostas de cobertura), em que licitantes submetem lances elevados para simular disputa e viabilizar a vitória de um parceiro~\cite{carbone2024bidrigging,aaltio2025complementary}; (ii) supressão de propostas, caracterizada pela abstenção coordenada de concorrentes aptos; (iii) \textit{bid rotation} (rodízio de vencedores), marcado pela alternância cíclica de vitórias ao longo de certames sucessivos~\cite{imhof2018screening,imhof2021detecting}; e (iv) alocação setorial ou territorial, na qual fornecedores dividem secretarias ou nichos específicos, elevando artificialmente a concentração orçamentária~\cite{wachs2019cartel,wachs2020corruption}. Como os participantes costumam cumprir formalmente os requisitos habilitatórios, a detecção preventiva depende de métricas objetivas e relacionais capazes de rastrear padrões anômalos em grandes volumes de registros~\cite{fazekas2016objective}.

\subsection{Modelagem em Redes Complexas para Auditoria Pública}

A Ciência das Redes permite traduzir bases administrativas relacionais em estruturas formais analisáveis algoritmicamente~\cite{newman2018networks}. Em contratações públicas municipais brasileiras, estudos recentes demonstram a eficácia de projeções empresa--empresa de co-disputa~\cite{lyra2021characterization} e de grafos heterogêneos que acoplam co-participação em certames a vínculos societários do QSA~\cite{villamil2024collusion,pompeu2025collusion}. Para capturar integralmente o ciclo das despesas municipais sem distorcer a semântica de cada etapa, este trabalho adota três representações complementares:

\begin{enumerate}
  \item \textbf{Grafo Heterogêneo Multicamada ($G_{\text{soc}}$):} definido como $G_{\text{soc}} = (V_{\text{soc}}, E_{\text{soc}}, \tau, \varphi, w)$, em que $\tau$ particiona $V_{\text{soc}}$ nas classes \textit{Órgão}, \textit{Empresa} e \textit{Sócio Real}, e $\varphi$ distingue arestas de participação, co-concorrência, propriedade societária (\textit{É\_Dono\_De}) e vínculo cadastral ou societário suspeito (\textit{Vinculo\_Suspeito}), permitindo arestas paralelas ponderadas por $w: E_{\text{soc}} \to \mathbb{R}_{>0}$.
  \item \textbf{Dígrafo Ponderado de Disputas ($G_{\text{dir}}$):} definido como $G_{\text{dir}} = (V_{\text{emp}}, E_{P \to V}, w_{\text{dir}})$, em que uma aresta direcionada $u \to v$ registra que a empresa perdedora $u$ foi superada pela vencedora $v$ em itens multi-participante. Em bases públicas como o SICOM~\cite{tcemg2025in01}, nas quais os lances perdedores frequentemente aparecem zerados nos registros de julgamento, a modelagem direcionada preserva a assimetria vitória/derrota e viabiliza a análise de reciprocidade temporal e supressão de desconto frente à cotação prévia~\cite{imhof2021detecting,wachs2019cartel}.
  \item \textbf{Grafo Bipartido Unidade--Credor ($G_{\text{bip}}$):} modelado como rede de dois modos $G_{\text{bip}} = (U \cup C, E_{\text{emp}}, w_{\text{fin}})$~\cite{borgatti1997network}, ligando unidades orçamentárias $u \in U$ aos credores $c \in C$ com peso igual ao valor líquido empenhado. Aplicações similares na América Latina evidenciam que a análise bipartida revela dependência crônica e fornecedores-ponte em administrações descentralizadas~\cite{herrera2020characterization,alcazar2026supplier}.
\end{enumerate}

\subsection{Algoritmos de Grafos e Métricas Aplicadas}

Para processar as três representações dentro de limites determinísticos de complexidade assintótica, emprega-se o seguinte conjunto de algoritmos e métricas:

\subsubsection{Redução por $1$-Core, Union-Find e Listagem de Triângulos $K_3$}
Dado o grafo heterogêneo $G_{\text{soc}}$, filtra-se o subgrafo induzido pelas arestas \textit{Vinculo\_Suspeito} acima do limiar de risco e suas respectivas arestas de propriedade. Em seguida, aplica-se a decomposição $k$-core com $k=1$~\cite{seidman1983network} via fila em busca em largura ($\mathcal{O}(|V| + |E|)$), podando vértices isolados. Sobre o subgrafo resultante, a estrutura \textit{Disjoint-Set Union} (Union-Find) com compressão de caminhos e união por posto~\cite{tarjan1975efficiency} extrai as componentes conexas em tempo quase linear $\mathcal{O}(|E|\,\alpha(|V|))$, onde $\alpha$ é a inversa da função de Ackermann. Por fim, para cada par suspeito $(a,b)$, interseccionam-se as listas de adjacência de sócios em tempo $\mathcal{O}(|E|\,d_{\max})$ para enumerar triângulos societários fechados ($K_3$: Sócio--Empresa$_A$--Empresa$_B$).

\subsubsection{Componentes Fortemente Conexas (Tarjan), Ciclos Curtos e Comunidades de Louvain}
No dígrafo $G_{\text{dir}}$, o algoritmo de Tarjan~\cite{tarjan1972depth} identifica todas as Componentes Fortemente Conexas (SCCs) em uma única travessia em profundidade com complexidade $\mathcal{O}(|V| + |E|)$. Enquanto relacionamentos acíclicos geram SCCs unitárias ($|C| = 1$), a alternância de vitórias entre concorrentes produz SCCs não triviais ($|C| \ge 2$). Para isolar alternâncias restritas a múltiplas licitações, executa-se uma busca em profundidade limitada para ciclos simples $C_3$ e $C_4$ ($\mathcal{O}((|V|+|E|)c)$, sendo $c$ o fator de ramificação local). Complementarmente, o algoritmo de Louvain~\cite{blondel2008fast} maximiza a modularidade $Q$ na projeção não direcionada de $G_{\text{dir}}$:
\begin{equation}
  Q = \frac{1}{2m} \sum_{x, y \in V_{\text{emp}}} \left[ A_{xy} - \frac{k_x k_y}{2m} \right] \delta(g_x, g_y),
\end{equation}
onde $A_{xy}$ é o peso da aresta, $k_x = \sum_y A_{xy}$, $m = \frac{1}{2}\sum_{xy} A_{xy}$, $g_x$ é a comunidade do vértice $x$ e $\delta$ é a função delta de Kronecker. Ressalta-se que valores elevados de $Q$ e a presença de ciclos $C_3/C_4$ indicam apenas segmentação de mercado ou disputa recorrente entre especialistas de um mesmo nicho, exigindo confronto com razões de preço e datas distintas de certames~\cite{imhof2018screening,aaltio2025complementary}. Para caracterização estrutural comparativa da rede de co-disputas no painel de síntese, utiliza-se também a centralidade de intermediação não ponderada $C_B(v)$~\cite{freeman1977set}, computada pelo algoritmo de Brandes~\cite{brandes2001faster} em tempo $\mathcal{O}(|V|\cdot|E|)$.

\subsubsection{Concentração Orçamentária (HHI, Entropia de Shannon e PageRank Bipartido)}
No grafo bipartido $G_{\text{bip}}$, para cada unidade orçamentária $u \in U$ com orçamento empenhado total $W_u = \sum_{c \in N(u)} w_{u,c}$ e participação relativa do credor $s_{u,c} = w_{u,c} / W_u$, calcula-se o Índice Herfindahl--Hirschman ($\text{HHI}_u$) em $\mathcal{O}(|E|)$~\cite{rhoades1993herfindahl}:
\begin{equation}
  \text{HHI}_u = 10.000 \sum_{c \in N(u)} s_{u,c}^2 \in (0,\, 10.000].
\end{equation}
Seguindo as faixas referenciais do DOJ/FTC~\cite{doj2010horizontal}, valores abaixo de $1.500$ indicam baixa concentração, de $1.500$ a $2.499$ concentração moderada e $\ge 2.500$ alta concentração. Para detectar dispersão versus dependência em poucos fornecedores, calcula-se a Entropia de Shannon Normalizada~\cite{shannon1948mathematical} sobre a distribuição orçamentária da unidade ($N_u = |N(u)| > 1$):
\begin{equation}
  H_{\text{norm}}(u) = \frac{-\sum_{c \in N(u)} s_{u,c} \ln(s_{u,c})}{\ln(N_u)} \in [0,\, 1],
\end{equation}
em que valores próximos de zero denotam extrema concentração na unidade $u$. Aplicando contração de vértices sobre credores ligados ao mesmo componente societário do Eixo 1, computa-se $\Delta\text{HHI}_u = \text{HHI}_{u,\text{grupo}} - \text{HHI}_u$. Por fim, o algoritmo \textit{PageRank} ponderado~\cite{brin1998anatomy} ($\mathcal{O}(|V|+|E|)$ por iteração) ranqueia a relevância sistêmica conjunta das secretarias concentradoras e dos credores que atuam como pontes entre múltiplas pastas municipais.


% 3. Metodologia (Dataset, Modelagem V e E, Algoritmos e Fluxograma Obrigatório)
% =============================================================================
% SEÇÃO 3: METODOLOGIA
% Diretrizes do Professor (Seções 3.2 e 3.3 do edital):
% - A "receita de bolo" técnica do seu projeto.
% - Origem e natureza dos dados utilizados (dataset, formato, pré-processamento).
% - Rigoroso processo de modelagem: como os elementos do problema real foram
%   traduzidos computacionalmente em vértices (V) e arestas (E) e/ou nós de uma árvore?
% - Descrição da lógica do algoritmo proposto passo a passo.
% - OBRIGATÓRIA A INCLUSÃO DE UMA FIGURA CONTENDO O FLUXOGRAMA DA SUA PROPOSTA,
%   mapeando o ciclo de vida da solução da entrada de dados até a obtenção dos resultados.
% =============================================================================

\section{Metodologia}
\label{sec:metodologia}

% TODO: Escreva a Metodologia aqui:
% 1. Origem e natureza da base de dados;
% 2. Processo formal de modelagem matemática (Definição de Vértices, Arestas e Pesos);
% 3. Lógica algorítmica passo a passo;
% 4. Inclusão e discussão do Fluxograma Obrigatório.


A metodologia deste estudo organiza-se em um pipeline estruturado em três eixos: (i) a ingestão, limpeza  e consolidação da base dados; (ii) modelagem relacional das vertentes de investigação e suas respcecitvas modelagens computacional baseadas em grafos; (iii) enquadramento e inferência dos resultados.


\subsection{Banco de Dados e Pipeline de Ingestão}

A obtenção de dados sobre compras públicas municipais no Brasil frequentemente esbarra na fragmentação dos portais locais. No município de Divinópolis/MG, o portal de transparência municipal~\cite{portal_transparencia_betha} fornece apenas consultas agregadas, impedindo o rastreamento atômico de lances e certames. Em contrapartida, o Portal de Dados Abertos do Tribunal de Contas do Estado de Minas Gerais (TCE-MG)~\cite{tce_mg_dados_abertos} centraliza as remessas formais obrigatórias do Sistema Informatizado de Contas dos Municípios (SICOM). Para este estudo, extraiu-se o conjunto completo de dados jurisdicionados brutos do município referentes ao exercício financeiro de 2025 integralmente consolidado.

Na sua forma bruta original, a base compreende 55 tabelas em formato CSV (delimitadas por ponto e vírgula), distribuídas em cinco módulos orçamentários: estrutura organizacional (\texttt{orgao}), certames e disputas (\texttt{licitacao}), instrumentos contratuais (\texttt{contrato}), reservas orçamentárias (\texttt{empenho}) e liquidações financeiras (\texttt{despesa}). A manipulação direta desses arquivos impunha alta redundância de metadados protocolares de envio e ruídos cadastrais (como perda de zeros à esquerda em documentos fiscais e valores monetários em texto regional).

Para viabilizar a modelagem computacional, desenvolveu-se um pipeline de Extração, Transformação e Carga (ETL) utilizando o motor analítico colunar \texttt{DuckDB} integrado a um banco relacional SQLite (\texttt{sicom.db}). A transformação guiou-se por três operações centrais: (i) projeção seletiva das entidades e tabelas relevantes para a rede (\texttt{secretarias}, \texttt{licitacoes}, \texttt{itens}, \texttt{propostas\_lances}, \texttt{homologacoes} e \texttt{empresas}); (ii) padronização vetorial via \texttt{LPAD} para recompor os 14 dígitos de CNPJs e 11 de CPFs, eliminando a criação espúria de vértices duplicados; e (iii) conversão de tipos de dados (datas em padrão ISO-8601 e valores monetários para ponto flutuante de dupla precisão \texttt{FLOAT64}). Esse saneamento reduziu em mais de 80\% o volume de dados redundantes, gerando uma base canônica e indexada para a construção dos grafos.

\subsection{Modelagem relacional e computacional dos grafos}


A investigação do ciclo de compras públicas municipais foi estruturada em três eixos analíticos complementares, cada um responsável por capturar uma dimensão específica do comportamento dos agentes econômicos e da administração pública. A partir do banco de dados canônico (sicom.db), cada eixo foi   formalmente traduzido em estruturas topológicas de grafos, definindo-se rigorosamente seus conjuntos de vértices, arestas, direcionamentos e funções de   ponderação

\subsubsection{}


% -----------------------------------------------------------------------------
% EXEMPLO DE INCLUSÃO DA FIGURA DO FLUXOGRAMA (OBRIGATÓRIO):
% -----------------------------------------------------------------------------
% \begin{figure}[!htbp]
% \centering
% \includegraphics[width=\columnwidth]{figures/fluxograma.png}
% \caption{Fluxograma metodológico do ciclo de vida da solução computacional proposta.}
% \label{fig:fluxograma}
% \end{figure}


% 4. Resultados e Discussão (Tabelas, Gráficos, Gargalos e Análise Crítica)
% =============================================================================
% SEÇÃO 4: RESULTADOS E DISCUSSÃO
% Diretrizes do Professor (Seção 3.2 do edital):
% - O clímax do trabalho.
% - Apresente gráficos de desempenho, tempo de execução ou tabelas de assertividade.
% - Discuta criticamente os resultados (mais do que apenas mostrar números):
%   * O que os resultados significam?
%   * Houve algum gargalo computacional?
%   * O algoritmo se comportou como esperado quando a entrada cresceu (complexidade)?
% =============================================================================

\section{Resultados e Discussão}
\label{sec:resultados}

% TODO: Escreva a análise de Resultados e Discussão aqui:
% 1. Apresentação das métricas topológicas e estatísticas computadas;
% 2. Tabelas comparativas de resultados e assertividade;
% 3. Gráficos de desempenho e visualizações das redes;
% 4. Discussão crítica aprofundada dos significados e anomalias encontradas;
% 5. Avaliação do comportamento assintótico e escalabilidade dos algoritmos.

% -----------------------------------------------------------------------------
% EXEMPLO DE INCLUSÃO DE TABELA:
% -----------------------------------------------------------------------------
% \begin{table}[!htbp]
% \centering
% \caption{Métricas estruturais das redes analisadas.}
% \label{tab:metricas}
% \begin{tabular}{lcc}
% \toprule
% \textbf{Métrica} & \textbf{Descrição} & \textbf{Valor} \\
% \midrule
% $|V|$ & Número de Vértices & -- \\
% $|E|$ & Número de Arestas & -- \\
% \bottomrule
% \end{tabular}
% \end{table}

% -----------------------------------------------------------------------------
% EXEMPLO DE INCLUSÃO DE GRÁFICO / FIGURA:
% -----------------------------------------------------------------------------
% \begin{figure}[!htbp]
% \centering
% \includegraphics[width=\columnwidth]{figures/grafico_exemplo.png}
% \caption{Descrição do gráfico e comportamento observado.}
% \label{fig:grafico}
% \end{figure}


% 5. Conclusão e Próximas Etapas (Contribuições, Limitações e Trabalhos Futuros)
% =============================================================================
% SEÇÃO 5: CONCLUSÃO E PRÓXIMAS ETAPAS
% Diretrizes do Professor (Seção 3.2 do edital):
% - Fechamento conciso das contribuições da equipe.
% - Relato dos desafios e limitações práticas enfrentadas.
% - Sugestões de melhorias e direções para pesquisas futuras.
% =============================================================================

\section{Conclusão e Próximas Etapas}
\label{sec:conclusao}

% TODO: Escreva a Conclusão aqui:
% 1. Fechamento e principais contribuições alcançadas pelo projeto;
% 2. Principais desafios técnicos e limitações dos dados/modelagem;
% 3. Propostas de extensões e trabalhos futuros.


% =============================================================================
% REFERÊNCIAS BIBLIOGRÁFICAS (Arquivo references.bib)
% Regra do edital: Somente citações científicas reais (scholar.google.com).
% Deve conter link do repositório (GitHub) e base de dados (Zenodo).
% =============================================================================
\bibliographystyle{IEEEtran}
\bibliography{references}

\end{document}
